\scanpage{149}
\deleted{Indépendamment de ces faits, \inserted{et sans hyp. spécifique ``crois. normaux''} on voit que $\mathcal V^{*}_{D_1,D_0=X}$ \deleted{peut} s'interprète comme $\mathcal V_{\Theta^{(1)},\unsure{X}}\setminus\Theta^{(1)}$ --- il donne la ``structure à l'infini''}

Reprenons la situation du pt de vue topologique (sans nous \unsure{focuser} sur les $\Pi_1$). Le diviseur $\Theta$ à croisements normaux \inserted{sur $X$} définit une filtration décroissante du \deleted{$X$} analytico-topologique $X$
\[
(244)\qquad\Theta^{(0)}=X\supset\Theta^{(1)}=\Theta\supset\Theta^{(2)}\supset\Theta^{(\unsure{3})}\cdots
\]
On peut définir
\[
D_d=\widetilde{\Theta^{(d)}}\qquad\text{``normalisé topologique'' de }\Theta^{(d)},\text{ dans un sens à préciser }-----,
\]
mais cette définition dépend justement d'hyp ``crois. normaux'' \deleted{\unsure{des}} relativement délicates. Par contre, \deleted{\unsure{D}}
\[
\deleted{D_d^{*}=\Theta^{(d)}\setminus\Theta^{(d-1)}}
\]
\deleted{se définit sans telles hypothèses, donc aussi $\mathcal V_{D_d^{*}}$ et}
\[
\deleted{(245)}\qquad
\diagram{\begin{tikzcd}[ampersand replacement=\&]
\deleted{\mathcal V^{*}_{D_d^{*}}}\arrow[r,hook] \& \deleted{\mathcal V_{D_d^{*}}}
\end{tikzcd}}
\]
\margin{148}
